% !TEX root = z-main.tex

\section{Enfoque metodológico}
La investigación respondió a un enfoque aplicado, con características proyectivas, descriptivas y cualitativas, orientado al desarrollo de modelos tridimensionales de estructuras arqueológicas del sitio Kaminaljuyú. El proyecto buscó representar digitalmente estos elementos mediante técnicas de modelado \gls{3d}, basadas en datos arqueológicos previamente verificados y validados por profesionales en el área, sin manipulación de variables cuantitativas. Se trató de una investigación enfocada en la reconstrucción visual y contextual de un sitio patrimonial con fines educativos, culturales y tecnológicos.

El desarrollo se realizó en un entorno interdisciplinario que integró conocimientos de arqueología, diseño digital y tecnologías de visualización. La planificación y el seguimiento de las actividades se gestionaron por medio de una plataforma digital colaborativa, la cual permitió asignar tareas semanales, validar avances y centralizar la documentación del proceso.

La planificación y el seguimiento del proyecto se gestionaron mediante una plataforma digital colaborativa, Notion, donde el equipo encargado de la gestión asignó semanalmente las tareas de modelado, organizó las prioridades y centralizó los avances. Esta plataforma permitió documentar el flujo de trabajo, dar seguimiento individual al progreso de las tareas y garantizar una trazabilidad clara del proceso. El modelado de los montículos se basó en información técnica proporcionada por especialistas en arqueología, quienes entregaron croquis, medidas e imágenes de referencia que posibilitaron representar las estructuras con fidelidad.

Una vez recibida la información, se configuró el entorno de trabajo en el software de modelado Blender, donde se establecieron las escalas apropiadas y se organizaron carpetas por montículo. La construcción tridimensional de las estructuras se inició desde las bases y avanzó progresivamente hacia las secciones superiores. Se emplearon técnicas mixtas de modelado de caja, y se previó un ritmo de producción de aproximadamente un modelo por semana, ajustado a la complejidad de cada estructura y a la retroalimentación técnica y arqueológica recibida.

Las texturas y materiales utilizados en los modelos fueron seleccionados desde el repositorio de recursos digitales Poly Haven y se aplicaron bajo la supervisión del área de arqueología, garantizando coherencia visual y contextual con los materiales históricos originales, como el barro, la piedra y los techados vegetales. Se prestó especial atención a la correcta proyección UV, evitando estiramientos o errores de mapeo en las texturas aplicadas. Cada modelo completado fue validado por especialistas arqueológicos, quienes emitieron observaciones sobre proporciones, formas y texturas. En función de estas observaciones, se realizaron los ajustes necesarios para asegurar la fidelidad de las representaciones.

Como parte del proceso de mejora técnica, se previó una participación en la comunidad oficial de Blender. Se estableció contacto con expertos para solicitar retroalimentación sobre aspectos clave del modelado. Las preguntas que se realizaron incluyeron si la topología del modelo era adecuada para su uso en realidad aumentada, si las proporciones y la escala estaban correctamente definidas dentro del espacio tridimensional, y si las texturas habían sido aplicadas sin errores visibles de mapeo UV. Se esperó que la comunidad respondiera con recomendaciones puntuales y videos explicativos, los cuales fueron analizados y aplicados para mejorar los estándares técnicos del trabajo. Esta etapa complementó la validación arqueológica, fortaleciendo la solidez técnica de los modelos que se generaron.

Los modelos finalizados fueron almacenados en un repositorio privado en GitHub, en formatos .obj y .mtl, preparados para su integración posterior. Los archivos fuente en formato .blend se conservaron localmente y estuvieron respaldados automáticamente en la nube mediante OneDrive, garantizando así la seguridad y la integridad del material original. Todo el proceso fue documentado gráficamente mediante capturas de pantalla y registros de validación, centralizados en Notion y WhatsApp.

Aunque la integración de los modelos tridimensionales en la aplicación de realidad aumentada no formó parte directa de las tareas del área de modelado, se tuvo prevista la participación en dicha fase. Esta colaboración permitió verificar el comportamiento de los modelos en el entorno digital, identificar posibles ajustes estructurales o visuales y asegurar la compatibilidad técnica con motores gráficos como \gls{ARCore}.

El análisis de los modelos se llevó a cabo mediante un proceso de validación técnica y visual, sustentado en la retroalimentación de profesionales en arqueología, quienes evaluaron aspectos como proporción, escala y coherencia estética. Desde una perspectiva técnica, se revisó la topología, el mapeo de texturas y la compatibilidad con entornos digitales. Este análisis fue de tipo cualitativo, iterativo y no estadístico, dado que se orientó a la mejora continua de los modelos a través de observación crítica y ajuste progresivo.

Finalmente, se observaron consideraciones éticas relacionadas con la representación y difusión del patrimonio cultural. Este proyecto no implicó la recolección de datos personales ni el involucramiento de individuos o comunidades. Sin embargo, se asumió el compromiso de representar con precisión los elementos arqueológicos modelados, evitando distorsiones que pudieran inducir a errores históricos o técnicos. Se mantuvo una actitud de respeto hacia la memoria cultural del pueblo maya, utilizando las tecnologías digitales con fines educativos, académicos y de conservación.

\section{Etapas de desarrollo}
\subsection{Planificación y asignación de tareas}
La gestión del proyecto se llevó a cabo a través de la plataforma Notion, en la cual el equipo responsable de la coordinación general asignó semanalmente las tareas correspondientes a cada fase del modelado. Cada semana se estableció la elaboración de un montículo específico, en función de la disponibilidad de los datos técnicos proporcionados por el área de arqueología. El avance de cada actividad quedó documentado en dicha plataforma, lo que permitió su trazabilidad y el control de entregas.

\subsection{Recolección y análisis de información arqueológica}
El equipo arqueológico especializado proporcionó las medidas, croquis e imágenes de referencia necesarios para modelar cada una de las estructuras seleccionadas. Esta información fue analizada e interpretada para ser adaptada al entorno digital, asegurando la fidelidad estructural y contextual de los modelos.

\subsection{Configuración del entorno de trabajo en Blender}
Se utilizó el software Blender, versión 4.3.2, como herramienta principal para el modelado tridimensional. Se crearon entornos de trabajo independientes para cada estructura, organizados en carpetas con nombres específicos por montículo. Se verificó la escala general del espacio tridimensional con base en los datos arqueológicos recibidos, garantizando la coherencia entre los modelos digitales y las dimensiones reales.

\subsection{Modelado tridimensional de estructuras}
El proceso de construcción de los modelos se inició desde las bases estructurales hasta los niveles superiores, respetando las proporciones y características definidas en los registros arqueológicos. El ritmo de trabajo estimado fue de un modelo por semana, dependiendo de la complejidad de cada montículo y con una semana adicional destinada a la retroalimentación obtenida tras cada revisión.

\subsection{Aplicación de texturas}
Las texturas fueron seleccionadas desde el repositorio libre Poly Haven (https://polyhaven.com), con el acompañamiento del equipo arqueológico para asegurar coherencia con los materiales originales utilizados históricamente (barro, piedra, techados vegetales, entre otros). Se revisó la correcta proyección de las texturas para evitar estiramientos o distorsiones, realizando los ajustes necesarios cuando fue requerido.

\subsection{Validación arquelógica y ajustes}
Cada modelo finalizado fue enviado al equipo especializado en arqueología para su revisión técnica y contextual. Las observaciones emitidas fueron documentadas y, en caso de ser necesario, se implementaron los ajustes técnicos correspondientes en el diseño tridimensional.

\subsection{Asesoramiento técnico externo}
Durante el desarrollo del proyecto, se previó la participación en la comunidad oficial de Blender, donde se consultaron dudas específicas relacionadas con herramientas o procesos técnicos de modelado. Estas interacciones permitieron optimizar el flujo de trabajo y mejorar la calidad de los modelos generados.

\subsection{Evaluar la calidad y precisión del modelado 3d en su integración con la aplicación de realidad aumentada}
Se llevó a cabo la evaluación de la integración de los modelos \gls{3d} en la aplicación de realidad aumentada, verificando su precisión, calidad visual y comportamiento interactivo, e incorporando la retroalimentación de usuarios especializados en arqueología y tecnología para garantizar su fidelidad y compatibilidad técnica.

\subsection{Registro}
Durante todo el desarrollo se generó documentación gráfica mediante capturas de pantalla, así como versiones intermedias y finales de cada modelo compartidas por medio de WhatsApp. Las observaciones técnicas también fueron archivadas. Los archivos finales en formato .obj y .mtl fueron subidos a un repositorio privado en GitHub con fines de respaldo y distribución técnica. Los archivos fuente en formato .blend fueron almacenados localmente y respaldados de forma automática en la nube (OneDrive), evitando así pérdidas o ediciones no autorizadas.
